We present \aris{} (\textbf{A}utonomous \textbf{R}endering with \textbf{I}terative \textbf{S}cene-aware feedback),
a fully automated pipeline for constructing training-ready synthetic datasets for image editing tasks.
Building high-quality paired training data for 3D-aware image editing---such as view-controlled rotation---
has traditionally required extensive manual tuning of rendering parameters in tools like Blender,
an effort that is both time-consuming and difficult to scale.
\aris{} addresses this by introducing a \emph{free-form VLM-guided rendering evolution loop}:
a vision-language model reviews rendered scenes using natural language feedback,
and an AI agent reads this feedback directly to select the next rendering action
from a structured action space, iteratively improving visual quality
until the reviewer confirms the result is acceptable.
We instantiate this pipeline to construct \textbf{ARIS-Rotate},
a 50-object, 350-pair rotation editing dataset in which each sample pairs
a canonical front-view image with a target view specified by a natural-language prompt
(e.g., \textit{``Rotate this object from front view to right side view''}).
Ablation experiments demonstrate that the VLM evolution loop substantially improves
rendering quality over single-pass rendering without iterative refinement.
Fine-tuning Qwen Image Edit 2511~\cite{PLACEHOLDER_qwen_image_edit_verify} with LoRA on ARIS-Rotate
yields measurable improvements over the base model in PSNR, SSIM, and LPIPS on a held-out test set,
validating the downstream training value of automatically constructed data.
